% Source: Weinberg GR, physical PDF pp. 57--58
%         (printed pp. 31--32).
% Coverage: Section 2.3, Particle Dynamics; equations (2.3.1)--(2.3.8).
% Figures/tables/footnotes: none.
% Status: source-reviewed and compile-clean.
% Uncertainties: none.

\subsection{Particle Dynamics}\label{sec:2.3}

Let us suppose that a particle moves in a field of force at a velocity so high
that Newtonian mechanics does not suffice to calculate its motion. Let us also
suppose, as in the case of electrodynamics, that we know how to calculate the
force \(\mathbf F\) on our particle in any Lorentz frame in which, at a given
moment, it is at rest. Then we could compute the motion of our particle by
performing a Lorentz transformation to a frame in which the particle is at rest
at some time \(t_0\), computing the velocity
\(d\mathbf v=\mathbf F\,dt/m\) at the time \(t_0+dt\), performing another
Lorentz transformation to bring the velocity to zero again, and so on.
Fortunately, there is an easier way.

Let us define the \emph{relativistic force} \(f^\mu\) acting on a particle with
coordinates \(x^\mu(\tau)\) by
\begin{equation}
  f^\mu=m\frac{d^2x^\mu}{d\tau^2}.
  \tag{2.3.1}\label{eq:2.3.1}
\end{equation}
Clearly, if \(f^\mu\) were known, we could compute the motion of our particle.
We shall relate \(f^\mu\) to the Newtonian force by noting two of its properties:

\noindent(A) If the particle is momentarily at rest, then the proper time
interval \(d\tau\) equals \(dt\), so \(f^i=F^i\), where \(F^i\) are the
Cartesian components of the nonrelativistic force \(\mathbf F\), and
\begin{equation}
  F^0=0.
  \tag{2.3.2}\label{eq:2.3.2}
\end{equation}

\noindent(B) Under a general Lorentz transformation \eqref{eq:2.1.1}, the
coordinate differentials transform according to
\(dx'^\mu=\Lambda^\mu{}_\nu dx^\nu\), while \(d\tau\) is invariant, so
Eq.~\eqref{eq:2.3.1} tells us that \(f^\mu\) has the Lorentz transformation
rule
\begin{equation}
  f'^\mu=\Lambda^\mu{}_\nu f^\nu .
  \tag{2.3.3}\label{eq:2.3.3}
\end{equation}
Any quantity such as \(dx^\mu\) or \(f^\mu\) that transforms according to
Eq.~\eqref{eq:2.3.3} is called a \emph{four-vector}.

Now suppose that our particle has velocity \(\mathbf v\) at some moment \(t_0\),
and introduce a new coordinate system \(x'^\mu\), defined by
\[
  x^\mu=\Lambda^\mu{}_\nu(\mathbf v)x'^\nu ,
\]
where \(\Lambda(\mathbf v)\) is the ``boost'' defined by
Eqs.~\eqref{eq:2.1.17}--\eqref{eq:2.1.21}. Since
\(\Lambda(\mathbf v)\) is constructed so as to carry a particle from rest to
velocity \(\mathbf v\), and since our particle has velocity \(\mathbf v\) at
time \(t_0\) in the coordinate system \(x^\mu\), it must be at rest at this
moment in the coordinate system \(x'^\mu\). Hence, according to (A), the force
four-vector in the coordinate system \(x'^\mu\) at time \(t_0\) is equal to the
nonrelativistic force \(F^\mu\). And therefore, according to (B), the force in
our original coordinate system is
\begin{equation}
  f^\mu=\Lambda^\mu{}_\nu(\mathbf v)F^\nu ,
  \tag{2.3.4}\label{eq:2.3.4}
\end{equation}
or more explicitly, since \(F^0=0\),
\begin{equation}
  \mathbf f=\mathbf F+(\gamma-1)\mathbf v\,
  \frac{\mathbf v\cdot\mathbf F}{\mathbf v^2},
  \tag{2.3.5}\label{eq:2.3.5}
\end{equation}
\begin{equation}
  f^0=\gamma\,\mathbf v\cdot\mathbf F=\mathbf v\cdot\mathbf f ,
  \tag{2.3.6}\label{eq:2.3.6}
\end{equation}
with \(\mathbf v\) the instantaneous velocity.

Now that we know how to calculate \(f^\mu\), we can use the differential
equations \eqref{eq:2.3.1} to calculate the four dependent variables
\(x^\mu(\tau)\), and then eliminate \(\tau\) to determine \(\mathbf{x}(t)\).
However, the initial values of \(dx^\mu/d\tau\) must be chosen so that
\(d\tau\) really is the proper time; that is, so that
\begin{equation}
  -1=\eta_{\mu\nu}
  \frac{dx^\mu}{d\tau}\frac{dx^\nu}{d\tau}.
  \tag{2.3.7}\label{eq:2.3.7}
\end{equation}
Note that Eq.~\eqref{eq:2.3.7} will be true for all \(\tau\) if it is true at
some initial \(\tau\), providing that its derivative vanishes; that is,
providing that
\begin{equation}
  0=2\eta_{\mu\nu}f^\mu\frac{dx^\nu}{d\tau}.
  \tag{2.3.8}\label{eq:2.3.8}
\end{equation}
That this is true can be seen either directly from Eq.~\eqref{eq:2.3.4}, or
more elegantly by noticing that the right-hand side is Lorentz invariant:
\[
\begin{aligned}
  \eta_{\mu\nu}f'^\mu\frac{dx'^\nu}{d\tau}
  &=\eta_{\mu\nu}\Lambda^\mu{}_\rho\Lambda^\nu{}_\sigma
    f^\rho\frac{dx^\sigma}{d\tau}\\
  &=\eta_{\rho\sigma}f^\rho\frac{dx^\sigma}{d\tau},
\end{aligned}
\]
and that it vanishes by virtue of Eq.~\eqref{eq:2.3.2} in a reference frame in
which the particle is at rest.
